% 63e Olympiade Internationale de Mathématiques — Oslo, Norvège.
% Énoncés officiels en français (imo-official.org). Chaque problème vaut 7 points.

\begin{center}
  \emph{Premier jour --- lundi 11 juillet 2022 --- Durée : 4 heures 30}
\end{center}

\subsection*{Problème 1}
La banque d'Oslo émet des pièces d'argent (que l'on note $A$) et de bronze (que l'on note $B$). Morgane
possède $n$ pièces d'argent et $n$ pièces de bronze. Ces $2n$ pièces sont alignées de gauche à droite dans un
ordre arbitraire. On dit que des pièces consécutives faites du même métal forment une \emph{chaîne}. Étant
donné un entier $k$ fixé tel que $1 \leqslant k \leqslant 2n$, Morgane effectue plusieurs fois de suite
l'opération suivante : elle identifie la plus grande chaîne qui contient la $k$\textsuperscript{e} pièce en
partant de la gauche, puis elle déplace toutes les pièces de cette chaîne à gauche de toutes les autres pièces.
Par exemple, si $n = 4$ et $k = 4$, le processus partant de la configuration $AABBBABA$ sera
\[
AABBBABA \to BBBAAABA \to AAABBBBA \to BBBBAAAA \to BBBBAAAA \to \cdots.
\]
Trouver tous les couples $(n, k)$ tels que $1 \leqslant k \leqslant 2n$ et pour lesquels, quel que soit
l'ordre de départ, à un moment du processus, les $n$ pièces de gauche seront toutes formées du même métal.

\subsection*{Problème 2}
On note $\mathbb{R}_{>0}$ l'ensemble des réels strictement positifs. Trouver toutes les fonctions
$f \colon \mathbb{R}_{>0} \to \mathbb{R}_{>0}$ telles que, pour tout réel $x \in \mathbb{R}_{>0}$, il existe
exactement un réel $y \in \mathbb{R}_{>0}$ tel que
\[
x f(y) + y f(x) \leqslant 2.
\]

\subsection*{Problème 3}
Soit $k$ un entier naturel non nul et $S$ un ensemble fini de nombres premiers impairs. Démontrer qu'il
existe au plus une manière (à rotation et symétrie axiale près) de disposer les éléments de $S$ sur un
cercle de sorte que le produit de deux nombres placés l'un à côté de l'autre soit toujours de la forme
$x^2 + x + k$, avec $x$ entier naturel non nul.

\bigskip
\begin{center}
  \emph{Second jour --- mardi 12 juillet 2022 --- Durée : 4 heures 30}
\end{center}

\subsection*{Problème 4}
Soit $ABCDE$ un pentagone convexe tel que $BC = DE$. On suppose qu'il existe un point $T$ situé à
l'intérieur de $ABCDE$ tel que $TB = TD$, $TC = TE$ et $\widehat{ABT} = \widehat{TEA}$. On note $P$ et $Q$ les
points d'intersection respectifs des droites $(CD)$ et $(CT)$ avec la droite $(AB)$ ; on suppose que les
points $P$, $B$, $A$ et $Q$ sont alignés dans cet ordre. De même, on note $R$ et $S$ les points
d'intersection respectifs des droites $(CD)$ et $(DT)$ avec la droite $(AE)$, et on suppose que les points
$R$, $E$, $A$ et $S$ sont eux aussi alignés dans cet ordre.

Démontrer que les points $P$, $S$, $Q$ et $R$ sont cocycliques.

\subsection*{Problème 5}
Trouver tous les triplets $(a, b, p)$ d'entiers naturels non nuls tels que $p$ soit premier et
\[
a^p = b! + p.
\]

\subsection*{Problème 6}
Soit $n$ un entier naturel non nul. Un \emph{carré nordique} est une grille $n \times n$ dont les cases ont
été numérotées de $1$ à $n^2$ dans un ordre arbitraire, chaque numéro étant utilisé une seule fois. Deux
cases distinctes sont considérées comme adjacentes si elles ont un côté commun. Toute case dont le numéro
est inférieur aux numéros des cases qui lui sont adjacentes est appelée \emph{vallée}.
Un \emph{chemin de randonnée} est une suite formée d'au moins une case, telle que :
\begin{enumerate}
  \item[(i)] la première case de la suite est une vallée ;
  \item[(ii)] deux cases consécutives de la suite sont toujours adjacentes ;
  \item[(iii)] les cases de la suite sont numérotées dans l'ordre strictement croissant.
\end{enumerate}
Trouver, en fonction de $n$, le plus petit nombre possible de chemins de randonnée que peut contenir un
carré nordique.
